\documentclass[letterpaper,11pt]{article}
\usepackage[top=0.8in, bottom=0.8in, left=0.9in, right=0.9in]{geometry}
\usepackage[hidelinks]{hyperref}
\usepackage{lmodern}
\usepackage{parskip}
\usepackage{microtype}
\setlength{\parskip}{6pt}
\setlength{\parindent}{0pt}
\begin{document}
\begin{flushleft}
\textbf{\large Ajay Venkatesh} \\
Brooklyn, NY \\
+1 516 585 9013 \\
\href{mailto:av3855@nyu.edu}{av3855@nyu.edu}
\end{flushleft}
\vspace{10pt}
May 21, 2026
\vspace{6pt}

Hiring Manager \\
CGI

\vspace{10pt}

Dear Hiring Manager,

My name is Ajay Venkatesh. I'm finishing my M.S. in Computer Engineering at NYU, with production software experience at PwC in Bengaluru, a summer SDE internship at Amazon, and ongoing on-campus work at NYU's Novel AI Technologies. I'm writing about the Python Backend Engineer role at CGI.

\textbf{Python} is my daily driver for both backend services and ML pipelines. At Amazon I provisioned \textbf{Infrastructure-as-Code} on \textbf{AWS CDK} (Lambda, API Gateway, DynamoDB, IAM, CloudWatch) with multi-stage CI/CD and rollback controls, plus a high-throughput, event-driven processing system on \textbf{AWS ECS Fargate, SQS, SNS} that decoupled distributed message handling under variable load. At Campus Mesh I built scalable \textbf{Node.js microservices} powering real-time academic workflows backed by \textbf{MongoDB}, plus an LLM-powered AI assistant with \textbf{prompt engineering} and \textbf{RAG retrieval} over structured data.

On backend architecture: I'm comfortable designing REST APIs, working across relational and non-relational stores (\textbf{PostgreSQL}, \textbf{MySQL}, \textbf{MongoDB}, \textbf{DynamoDB}), and building production systems with monitoring, observability, and reliability under load. At PwC I owned a distributed backend on \textbf{Microsoft Azure} with event-driven processing, fault-tolerant retries, and parallelized batch ingestion that materially cut execution time on operational feeds.

On engineering discipline: I ship under code-review discipline with unit and integration testing as part of the build. At PwC I led code reviews and mentored engineers on software design principles, and at Amazon I used \textbf{Mockito} for backend testing of Java services. The same posture applies to Python development with \textbf{pytest} and standard linting workflows.

Stack-wise: \textbf{Python}, \textbf{Java}, \textbf{TypeScript}, \textbf{SQL}, plus \textbf{AWS}, \textbf{Docker}, \textbf{Kubernetes}, and \textbf{CI/CD} workflows. Comfortable across cloud platforms, distributed systems, and the testing discipline that consulting-style delivery demands.

What pulls me toward CGI is the operational seriousness of consulting-style work, where engineering across client engagements has to be both production-quality and shippable on real timelines, with documentation and accountability that compound across teams.

I'd welcome the chance to discuss further. Thanks for considering my application.

\vspace{12pt}
Sincerely, \\
Ajay Venkatesh

\end{document}